\documentclass[11pt]{article}
\usepackage[T1]{fontenc}
\usepackage[utf8]{inputenc}
\usepackage{lmodern}
\input{glyphtounicode}
\pdfgentounicode=1
\usepackage{amsmath,amssymb,amsthm,geometry,hyperref}
\geometry{margin=1in}
\hypersetup{colorlinks=true,linkcolor=blue,citecolor=blue,urlcolor=blue,pdftitle={Nonexistence of (243,121,60) difference sets in C3 × C9 × C9},pdfauthor={P. Graham},pdfsubject={Computer-assisted nonexistence in one finite abelian group},pdfkeywords={difference set, multiplier, exhaustive search, Lean}}
\newtheorem{theorem}{Theorem}
\newtheorem{lemma}{Lemma}
\title{Nonexistence of $(243,121,60)$ difference sets\\ in $C_3$ $\times$ $C_9$ $\times$ $C_9$}
\author{P. Graham\\ \small\url{https://orcid.org/0009-0007-8493-1505}}
\date{October 21, 2026}
\begin{document}
\maketitle
\begin{abstract}
A $(v,k,\lambda)$ difference set in a finite group of order $v$ is a
$k$-element subset in which every nonidentity element occurs exactly
$\lambda$ times as an ordered difference. We prove, by exhaustive
computation, that no $(243,121,60)$ difference set exists in
$C_3\times C_9\times C_9$. This group-specific case appears as unresolved
in the 1997 table of Vera L\'opez and Garc\'ia S\'anchez and was singled
out by Gordon in 2019 as the ``Smallest Open Case.'' A multiplier-fixed
translate has thirteen points in the subgroup of elements annihilated
by three. Moment identities, singleton congruences, and legal affine
normalizations reduce the problem to eight quotient multiplicity arrays,
each with a complete ternary lifting tree. Integer difference-count
searches exhaust these trees, with separate checks of decomposition
completeness. Lean proves the multiplier reduction, normalization,
lifting coverage, and pruning lemmas; its final nonexistence theorem is
explicitly conditional on the externally established search assertions.
\end{abstract}

The proofs, the Lean formalization, the search code, the certificate checkers and the text were produced by AI models (from OpenAI and Anthropic), working under the author's direction.

\section{The result and its historical setting}

Write $C_m=\mathbb Z/m\mathbb Z$ additively. For a subset $P$ of an
abelian group, put
\[
 d_P(g)=\bigl|\{(a,b)\in P\times P:a-b=g\}\bigr|.
\]
A $(v,k,\lambda)$ difference set in a group of order $v$ is a subset
$D$ with $|D|=k$ and $d_D(g)=\lambda$ for every $g\ne0$.

\begin{theorem}\label{thm:main}
There is no subset $D\subseteq C_3\times C_9\times C_9$ such that
$|D|=121$ and
\[
 \bigl|\{(a,b)\in D\times D:a-b=g\}\bigr|=60
 \qquad\text{for every }g\ne0.
\]
\end{theorem}

The theorem is computer-assisted. Its scope is exactly these parameters
in this group. The mathematical reduction and its coverage are proved
in Lean; the exhaustive integer computations that reject the covered
leaves are external to Lean. Section~\ref{sec:formal} states this
boundary precisely.

Vera L\'opez and Garc\'ia S\'anchez undertook a systematic study of
abelian difference sets with $100<k\leq150$~\cite{VeraGarcia1997}.
Their Table~1, printed page~105, marks the present group and parameters
with a question mark, meaning unknown according to their convention on
page~101. That paper was published in 1997; the publisher gives
February~28, 1997. This supplies a published unresolved-status record,
without asserting that it was the first occasion on which the question
was asked.

Gordon's La Jolla Difference Set Repository has made such group-specific
questions accessible and maintained their connection with the
literature. In his presentation of August~3, 2019, Gordon singled out
the present case as the ``Smallest Open Case''~\cite[slide~38]{Gordon2019}.
Here slide~38 is physical PDF page~49. The repository's archived
version~1.3, dated April~24, 2026, also records the exact entry
$\operatorname{DS}(243,121,60,[3,9,9])$ as Open~\cite{GordonData}.
The entry has a status and an empty comment, with no per-entry date.
The archive date therefore does not date the original insertion or
last review of this entry. Similarly, the website's September~3, 2026
update label is a site-level date. The entry's original insertion and
last-review dates have not been recovered.

These sources also put the question in perspective. For the same
parameters, the other six abelian groups of order $243$ are already
recorded as follows; this is historical context, not a new proof of
those entries:
\[
\begin{array}{c|c|l}
\text{group}&\text{recorded answer}&\text{attribution}\\\hline
C_{243}&\text{No}&\text{Lander}\\
C_3\times C_{81}&\text{No}&\text{Lander}\\
C_9\times C_{27}&\text{No}&\text{Arasu--Ma}\\
C_3^2\times C_{27}&\text{No}&\text{Vera L\'opez--Garc\'ia S\'anchez}\\
C_3^3\times C_9&\text{No}&\text{Vera L\'opez--Garc\'ia S\'anchez}\\
C_3^5&\text{Yes}&\text{Paley's finite-field construction}
\end{array}
\]
The source for the complete list and its attributions is
Gordon~\cite{GordonData}.\footnote{This small status summary is selected
and reformatted from Daniel M. Gordon's repository data, copyright
2024 Daniel M. Gordon, under
\href{https://creativecommons.org/licenses/by/4.0/}{CC BY 4.0}.
Group notation and presentation are changed; the source status
values are unchanged. No endorsement is implied.}
The 1997 table independently records the
negative entries for $C_{243}$, $C_3\times C_{81}$,
$C_3^2\times C_{27}$, and $C_3^3\times C_9$; its method~21 uses
multiplier-orbit enumeration with multiplier~61. Arasu and
Ma~\cite{ArasuMa2001} subsequently resolved $C_9\times C_{27}$.
For $C_3^5$, Paley's quadratic-residue construction~\cite{Paley1933}
in the additive group of $\mathbb F_{243}$ gives the stated parameters.
The 1997 table cites Emma Lehmer~\cite{Lehmer1953} for its positive entry.
The nonexistence at issue here concerns
the remaining group $C_3\times C_9\times C_9$.

The work uses the classical multiplier and character methods developed
by Hall~\cite{Hall1956}, Menon, Turyn, and others, and the systematic multiplier searches
of Vera L\'opez and Garc\'ia S\'anchez. It also benefits from the later
nonexistence and field-descent theory of Arasu, Ma, Leung, Schmidt, and
Miller. These contributions supply both the historical question and
the algebraic framework for reducing it to a finite computation.
The contributors to Lean, mathlib, and the declared library dependencies
are also acknowledged for the formal-development infrastructure.

\section{Published restrictions and their hypotheses}\label{sec:context}

Set
\[
 v=243=3^5,\qquad k=121,\qquad\lambda=60,\qquad
 n=k-\lambda=61,\qquad\exp(G)=9.
\]
The parameter identity $k(k-1)=\lambda(v-1)$ holds. The odd-order
Bruck--Ryser--Chowla equation $61x^2-60y^2=z^2$ has the nonzero
solution $x=y=z=1$. These elementary necessary conditions do not
exclude the parameters.

Turyn's character restrictions~\cite{Turyn1965} allow this application:
$\operatorname{ord}_3(61)=1$ is odd. At the nontrivial character
orders $3$ and $9$, the prime $61$ is not self-conjugate. Indeed, its
powers modulo $3$ equal $1$, and modulo $9$ cycle through $7,4,1$,
never $-1$. Obstructions requiring a common prime divisor of $n$ and
$v$ have no such prime here; in particular that ingredient of the
Lander-type restrictions in Leung--Ma--Schmidt~\cite{LeungMaSchmidt2004}
is absent. That comparison uses the stated common-prime setting,
without claiming a full reading of that paper's proof.
The multiplier theorem does apply and provides symmetry,
as proved in Section~\ref{sec:multiplier}; that symmetry alone is
not a nonexistence result.

For the field-descent function of Leung and
Schmidt~\cite[Definition~2.6]{LeungSchmidt2005}, one obtains
$F(9,61)=3$. Here $\operatorname{ord}_3(61)=1$ and
$61\not\equiv1\pmod9$, so the exponent $b_3=1$ is permitted in
that definition. Taking their factor $A$ to be the entire group and
the complementary factor trivial, Theorem~4.3 gives
\[
 \exp(G)\leq243\sqrt{\frac{3^2}{4\cdot61\cdot\varphi(3)}}
 =\frac{729}{\sqrt{488}}\approx33.0003,
\]
which permits exponent~9. Their Theorem~A cannot instead split off a
$C_9$ factor: its coprimality condition would require
$\gcd(9,\exp(C_3\times C_9))=1$, whereas this gcd is~9.
Theorem~4.1 requires a cyclic Sylow subgroup, also absent here.

Arasu and Ma's 2001 nonexistence result~\cite{ArasuMa2001} resolves
$C_9\times C_{27}$. The relevant theorem, as explicitly restated and
extended by Miller~\cite[Theorems~1 and~3]{Miller2025}, assumes
\[
 G=C_{p^e}\times J,\qquad e\geq2,\qquad
 \exp(\operatorname{Syl}_p J)<p^e.
\]
Splitting off $C_9$ in the present group leaves another $C_9$, so
the strict exponent inequality fails. The forbidden-multiplier
congruences would demand both $t\equiv4\pmod9$ and
$t\equiv1\pmod9$. Although the older numerical condition
$4(p-1)n=488>243$ holds, it cannot repair the failed group
hypothesis. This comparison uses Miller's restatement; the full
2001 original was not available for a complete reading.

Chen, Xiang, and Sehgal~\cite{ChenXiangSehgal1994} study skew Hadamard
abelian difference sets, which satisfy the additional partition
$G=\{0\}\mathbin{\dot\cup}D\mathbin{\dot\cup}(-D)$.
The present parameters do not supply this hypothesis. Consequently
a skew-only exponent bound cannot be applied to the ordinary
difference-set question. This scope distinction follows from the
published abstract; the full original was not obtained, and no
numerical specialization of its bound is needed here.

Recent work on signed difference sets by He and
Wu~\cite{HeWu2026} includes ordinary difference sets but does not
yield this exclusion by the criteria considered here. Their
Theorem~3.2 has the same failed self-conjugacy test for~61;
Theorem~4.2 requires even $v$; Theorem~4.4 requires $k=v-1$;
and Theorem~4.7 requires an elementary abelian group and a defect
bound, whereas this group is not elementary abelian and the defect
is $d=122>5$. Miller's and
He--Wu's cited versions are preprints, not claimed here as verified
journal publications.

These are explicit applications and failures of hypotheses, not a
claim to have exhausted all uses of these methods. The historical
Open entries document the question; they cannot establish an
exhaustive priority claim. Some older originals, including Lander's
Theorem~4.38~\cite{Lander1983} and the Hall, Paley, and Lehmer works
cited for historical credit, were not read in full. Hall's survey is
not used to infer that a cyclic-group statement alone covers every
abelian hypothesis of the multiplier specialization.
Theorem~\ref{thm:main} rests on the
reduction and exhaustion given here, rather than on an inference
from an Open label or from the absence of a located earlier answer.

\section{Multiplier fixedness and quotient moments}\label{sec:multiplier}

Use coordinates $G=C_3\times C_9\times C_9$ and put
\[
\begin{split}
 H&=G[3]=\{(x,3u,3w):x,u,w\in C_3\},\\
 K&=3G=\{(0,3u,3w):u,w\in C_3\}.
\end{split}
\]
Thus $|H|=27$, $|K|=9$, and
$Q=G/K\cong C_3^3$. In particular, $H$ and $K$ are different
subgroups. Elements of $G$ are encoded by $81x+9y+z$, using the
least nonnegative coordinate representatives; elements of $Q$ by
$9x+3\bar y+\bar z$.

\begin{lemma}\label{lem:fixed}
If $D$ is a difference set as in Theorem~\ref{thm:main}, then a
translate $D_0$ satisfies $61D_0=D_0$.
\end{lemma}
\begin{proof}
The first multiplier theorem applies because $G$ is abelian,
$n=61$ is prime, $61>\lambda=60$, and $\gcd(61,243)=1$.
It gives $61D=D+g$ for some $g\in G$. This is also the
$n_1=61$ specialization of Menon's multiplier theorem in
Miller~\cite[Theorem~5]{Miller2025}; see the classical character
and multiplier theory in Turyn~\cite{Turyn1965}.
Put $s=\sum_{d\in D}d$. Summing the set identity gives
$60s=121g$. Multiplication by $121$ has inverse multiplication
by~7 on $G$, since $121\cdot7\equiv1\pmod9$.
For $u=7s$ we therefore have $60u=g$, whence
\[
 61(D-u)=D+g-61u=D-u.
\]
Take $D_0=D-u$.
\end{proof}

The initial translating element $u$ can be anywhere in $G$.
Subsequent normalizing translations will lie in $H$.
Multiplication by $61$, equivalently by $7$ modulo $9$, has order
three and fixed subgroup exactly $H$. Its scalar orbits consist of
27 singletons and 72 triples. Every fixed subset is a union of
these orbits. A group automorphism commutes with scalar
multiplication, and translation by $t\in H$ preserves fixedness
because $60t=0$.

A difference set has no nonzero translation period: if $D+h=D$,
then $d_D(h)=121$, contrary to $60$ when $h\ne0$.
It follows that the fixed translates of a fixed difference set
$D_0$ are exactly $D_0+H$: fixedness of $D_0+t$ implies that
$60t$ is a period, hence $60t=0$, whose solution set in $G$ is $H$.

For later use, let $U\leq G$ and
$r_U(a+U)=|D\cap(a+U)|$. Counting ordered pairs whose
difference lies in $U$ gives
\begin{equation}\label{eq:moment}
 \sum_{a+U}r_U(a+U)=121,\qquad
 \sum_{a+U}r_U(a+U)^2=121+60(|U|-1).
\end{equation}
More generally, projected autocorrelation equals
$121+60(|U|-1)$ at zero and $60|U|$ elsewhere.
For an index-three subgroup, the three coset counts consequently
satisfy
\begin{equation}\label{eq:triples}
 t_0+t_1+t_2=121,\qquad
 t_0^2+t_1^2+t_2^2=4921,\qquad0\leq t_i\leq81.
\end{equation}
The complete integer solutions are the six permutations of
$(36,40,45)$. This small finite assertion is checked both in
the finite enumerations and in Lean.

\section{The thirteen singleton points}\label{sec:singletons}

Assume from now on that $D_0$ is fixed, and put $S=D_0\cap H$.
There are four projective nonzero linear forms on
$(\bar y,\bar z)\in C_3^2$. For each, the two nonzero fibers
contain only scalar triples, so their intersection counts with
$D_0$ are divisible by three. Of $36,40,45$, only $40$ is not
divisible by three; thus the zero-fiber count is~40.
Every point of $H$ lies in all four zero fibers, and every point
outside $H$ lies in exactly one. Hence
\[
 4\cdot40=4|S|+(121-|S|),\qquad |S|=13.
\]

Let $a,b,c$ be the sizes of $S$ in the three $x$-layers of $H$.
They satisfy $a+b+c=13$ and $0\leq a,b,c\leq9$.
On $Q=G/K$, the three special cells $(x,0,0)$ contain only
singleton points and have multiplicities $a,b,c$. Every other
cell contains three scalar triples, so its multiplicity lies
in $\{0,3,6,9\}$. The $U=K$ instance of
\eqref{eq:moment} gives
\[
 a^2+b^2+c^2\equiv601\equiv7\pmod9.
\]
The $x$-quotient in \eqref{eq:triples} forces the residues of
$(a,b,c)$ modulo three to be a permutation of $(1,0,0)$.
The complete list of layer triples satisfying these conditions
consists of the permutations of $(4,0,9)$ and $(4,3,6)$.

For any $h\in H\setminus\{0\}$, the scalar action permutes the
60 ordered pairs in $D_0^2$ of difference $h$. Its fixed pairs
have both endpoints in $S$; all remaining pair orbits have
length three. Therefore
\begin{equation}\label{eq:congruence}
 |S\cap(S+h)|\equiv0\pmod3
 \qquad(h\in H\setminus\{0\}).
\end{equation}
A $(4,0,9)$ layer pattern has intersection four with its
translate by $(1,0,0)$, so cannot satisfy \eqref{eq:congruence}.

For an explicit finite encoding, the point $(x,3u,3w)\in H$
has mask bit $9x+3u+w$, with bit zero the least significant bit.
The layer conditions above describe exactly
\begin{equation}\label{eq:maskcount}
 6\left[\binom94+\binom94\binom93\binom96\right]
 =5{,}335{,}092
\end{equation}
distinct subsets of $H$. This is a count of candidate singleton
subsets, not of difference sets or of all fixed 121-point sets.
Exhaustive testing of all 26 congruences leaves exactly 7,776
masks. An enumeration scanning every one of the
$\binom{27}{13}=20{,}058{,}300$ thirteen-point subsets
independently reproduces both counts.

The surviving set of masks is exactly the orbit of
\begin{equation}\label{eq:mask}
 S_0:\quad\text{mask }24909666,
 \qquad (a,b,c)=(4,3,6),
\end{equation}
under the legal affine maps described next. Equality of these
two finite sets, with normalization witnesses, establishes
normalization; agreement of their cardinalities alone would
not do so. Thus every hypothetical difference set has a fixed
legal image with $H$-part exactly $S_0$.

\section{Legality of the normalizing symmetries}\label{sec:symmetry}

Use $F=A+t$ with $t=(t_0,3t_1,3t_2)\in H$, where
\begin{equation}\label{eq:affine}
\begin{split}
 F(x,y,z)=\bigl(&ax+by+cz+t_0\pmod3,\\
              &3dx+ey+fz+3t_1\pmod9,\\
              &3gx+hy+iz+3t_2\pmod9\bigr).
\end{split}
\end{equation}
Here $a\ne0\pmod3$ and
$B=\left(\begin{smallmatrix}e&f\\h&i\end{smallmatrix}\right)$
is invertible modulo three. The parameters in the first
coordinate and $d,g,t_j$ are taken modulo three; entries of
$B$ have chosen lifts modulo nine.

The coordinate orders make $A$ additive. Its induced map on
$G/3G$ has block upper-triangular matrix with diagonal blocks
$a$ and $B\pmod3$, so is invertible. If $A(p)=0$, it follows
that $p\in3G=K$. On $K$, the map is again given by
$B\pmod3$, which is invertible; hence $p=0$.
Thus $A$ is an automorphism. Translation preserves differences,
and $60t=0$ gives $F(61p)=61F(p)$. Consequently $F$ preserves
the difference-set property and scalar fixedness. The affine
map $F$ need not fix zero; it is its linear part $A$ that is
a group automorphism.

For $b=c=0$, the induced $H$-actions, writing its last two
coordinates as a vector $\boldsymbol u\in C_3^2$, are
\[
 (x,\boldsymbol u)\longmapsto
 (ax+t_0,B\boldsymbol u+\boldsymbol d x+\boldsymbol t).
\]
They form a group of order
$2\cdot3\cdot48\cdot9\cdot9=23{,}328$.
Each has an explicit full-$G$ lift of \eqref{eq:affine}.
All these lifts have been checked on the full 243-point group
for bijectivity, translation in $H$, scalar commutation, and
additivity of the linear part on every generator edge.
The generator-edge identities imply additivity for arbitrary
pairs by telescoping along a coordinate-generator path.
The orbit of $S_0$ has size 7,776, so its stabilizer in this
$H$-action group has order three.

Since the three layer sizes of $S_0$ are distinct, an
$S_0$-stabilizing $H$-action has $a=1,t_0=0$. Each of its
three possibilities has nine $(b,c)$ extensions. These induce
27 distinct permutations on $Q$, forming a group of order~27.
Their chosen full-$G$ lifts preserve $S_0$, commute with the
scalar, and act consistently on every $K$-coset. These facts
are checked for the actual point maps, including full
pairwise additivity in the separate quotient-map check.
The chosen lifts need not themselves form a 27-element
subgroup of affine maps on $G$. The stabilizer of order three
on $H$ and the induced group of order~27 on $Q$ refer to
different actions; neither is asserted to be the full affine
stabilizer of $S_0$ on $G$.

\section{Complete quotient multiplicity classification}\label{sec:quotient}

Normalize $S$ to $S_0$ and define
\[
 r(q)=|D_0\cap(q+K)|\qquad(q\in Q).
\]
The notation now refers to the normalized fixed image of $D_0$.
This is an integer multiplicity array on $G/K$, not a subset
of $G$ and not an array on $G/H$. Its special entries
$r(x,0,0)$ are $(4,3,6)$; initially each of the other 24
entries may be any of $0,3,6,9$.

Choose one representative of each of the 13 projective
nonzero forms $L:Q\to C_3$, by requiring its first nonzero
coefficient to be~1, and put
\[
 t_L(j)=\sum_{L(q)=j}r(q).
\]
If $L$ ignores $x$, its zero fiber contains all thirteen
singleton points; otherwise it contains the four points of
the $x=0$ singleton layer. Thus $t_L(0)\equiv1\pmod3$ in
both cases. Equation~\eqref{eq:triples} forces
\begin{equation}\label{eq:lines}
 t_L(0)=40,\qquad \{t_L(1),t_L(2)\}=\{36,45\}.
\end{equation}

For $p,q\in Q$, exactly thirteen of the forms have
$L(p)=L(q)$ when $p=q$, and exactly four when $p\ne q$.
Therefore
\begin{equation}\label{eq:inversion}
 \sum_Lt_L(L(q))=13r(q)+4(121-r(q))=484+9r(q).
\end{equation}
This identity proves completeness of the following finite
procedure. Enumerate all $2^{13}=8192$ ways of ordering
$(36,45)$ in \eqref{eq:lines}; recover $r(q)$ from
\eqref{eq:inversion}; retain integer coefficients between
zero and nine, the special entries $(4,3,6)$, and nonspecial
entries divisible by three. Every possible quotient array
must occur, since its line assignments occur and the
inversion recovers it uniquely.

Exactly 144 arrays survive. Each has projected
autocorrelation 601 at zero and 540 at every nonzero
element of $Q$, as required by \eqref{eq:moment} and its
nonzero counterpart. Each has twelve nonspecial entries
equal to three and twelve equal to six. The absence of
zero and nine is an enumeration result, not an assumption
used to exclude possible arrays. A separate depth-first
enumeration assigns $\{0,3,6,9\}$ to the 24 nonspecial
cells and enforces all index-three moment triples,
without assuming which fiber has count~40. It gives
exactly the same set of 144 arrays.

The group of 27 induced $S_0$-stabilizing quotient
permutations partitions these arrays into eight orbits.
Representatives in the order used for lifting are given
in Table~\ref{tab:cases}. Each nine-digit string lists
$r(x,\bar y,\bar z)$ for fixed $x$, with
$3\bar y+\bar z=0,\ldots,8$ in increasing order.
\begin{table}[ht]
\centering
\begin{tabular}{c|r|ccc}
case $c$&orbit size&$x=0$&$x=1$&$x=2$\\\hline
0&27&\texttt{463666333}&\texttt{366636636}&\texttt{633336336}\\
1&27&\texttt{463663363}&\texttt{366633666}&\texttt{633333366}\\
2&9 &\texttt{436666333}&\texttt{366663663}&\texttt{633363363}\\
3&9 &\texttt{436636336}&\texttt{366633666}&\texttt{633333366}\\
4&27&\texttt{463666333}&\texttt{363633633}&\texttt{663366366}\\
5&27&\texttt{436666333}&\texttt{336633633}&\texttt{636366366}\\
6&9 &\texttt{463666333}&\texttt{366366663}&\texttt{633363633}\\
7&9 &\texttt{463636336}&\texttt{366336666}&\texttt{633333636}
\end{tabular}
\caption{The eight quotient multiplicity representatives.}
\label{tab:cases}
\end{table}
Their orbits are disjoint and their union is the full
144-array set. The orbit sizes are
$27,27,9,9,27,27,9,9$; their stabilizers in the induced
27-element group have orders $1,1,3,3,1,1,3,3$.
For each array an explicit legal map sends it to its
representative while preserving $S_0$. No search count is
divided by one of these symmetry orders.

\subsection{A centered coset invariant}\label{sec:invariant}

For the centered fixed translate, let
\[
 m(v)=|D_0\cap(v+H)|\qquad(v\in G/H\cong C_3^2).
\]
We have $m(0)=13$. The zero fiber for each of the four
directions on $G/H$ has count~40, so
\[
 m(v)+m(-v)=27\qquad(v\ne0).
\]
After singleton normalization, $m(v)=\sum_xr(x,v)$.
The complete quotient classification gives
$m(v)\in\{9,12,15,18\}$ for $v\ne0$.
Legal normalizations induce an invertible linear map on
$G/H$, so these statements transfer back to the centered
fixed translate. Each opposite pair is consequently
$\{9,18\}$ or $\{12,15\}$.

The $U=H$ moment identity is
\[
 \sum_vm(v)^2=121+60\cdot26=1681.
\]
If $j$ of the four opposite pairs are $\{9,18\}$, then
\[
 1681-13^2=j(9^2+18^2)+(4-j)(12^2+15^2)
 =1476+36j.
\]
Thus $j=1$: the eight off-$H$ cosets have multiplicities
\begin{equation}\label{eq:Hcosets}
 9,18,12,12,12,15,15,15,
\end{equation}
with one opposite pair $\{9,18\}$ and three
$\{12,15\}$. All 144 arrays satisfy this invariant.
It is a consequence of the classification, not an
additional unexplained pruning condition. For an
uncentered translate the distinguished thirteen-point
coset can move; $|D\cap H|=13$ is not asserted for an
arbitrary translate of the original difference set.

\section{Lifting completeness and the search partition}\label{sec:lifting}

Fix a case in Table~\ref{tab:cases}. Each of its 24
nonspecial $K$-cosets has nine points, partitioned into
three scalar triples. A multiplicity-three cell must
select exactly one triple, and a multiplicity-six cell
exactly two. In either case there are exactly three
choices, namely all members of the appropriate
$\binom31$ or $\binom32$ family.

Number a cell's triples $0,1,2$ by their least encoded
point. A digit $j$ in a multiplicity-three cell selects
triple $j$; in a multiplicity-six cell it omits triple
$j$. Order the cells by placing the twelve six-cells
first in increasing $Q$-index, followed by the twelve
three-cells in increasing $Q$-index. Denote the resulting
set for a case $c$ and digit word
$\boldsymbol d\in\{0,1,2\}^{24}$ by
$\mathcal L(c,\boldsymbol d)$.

\begin{lemma}\label{lem:lifting}
For a fixed case, the map $\boldsymbol d\mapsto
\mathcal L(c,\boldsymbol d)$ is a bijection onto all
scalar-fixed subsets with singleton part $S_0$ and that
quotient multiplicity array. Every such set has 121 points.
\end{lemma}
\begin{proof}
The union of $S_0$ and the chosen cell options has size
$13+12\cdot6+12\cdot3=121$, is scalar-fixed, and has
the prescribed multiplicities. Conversely, any fixed
set with those multiplicities intersects each cell
in an invariant subset of the required size, hence
in one of the three stated options. This proves
surjectivity. Intersecting the union with each of
the disjoint cells recovers its option and unique
digit, proving injectivity.
\end{proof}

There are $3^{24}=282{,}429{,}536{,}481$ syntactic
leaves per case before difference-count pruning.
The bijection does not require enumeration of all
these leaves. Different representative arrays give
disjoint families because their quotient arrays differ;
no assertion is made that all these leaves represent
distinct affine-equivalence classes.

The local families have been checked independently of
the search tables: scalar orbits were reconstructed
from modular coordinates and all their required
subfamilies formed, then compared as sets with every
option emitted by each compiled search kernel.
A separate completeness check also enumerated all
three-point and six-point subsets of each nine-point
cell and tested scalar closure directly. Both checks
cover all $8\cdot24\cdot3=576$ labeled options and
the thirteen initial points. The comparison definition
comes from scalar fixedness and subset size, rather
than from treating an existing option table as complete.

Split each 24-digit word after twelve digits. With
leading zeroes retained, give its first half the
big-endian rank
\[
 p=\sum_{j=0}^{11}d_j3^{11-j},\qquad
 \operatorname{id}=c\,3^{12}+p.
\]
The complete ID domain has size
\[
 N=8\cdot3^{12}=4{,}251{,}528,
\]
and is partitioned into 1,038 intervals
\begin{equation}\label{eq:ranges}
 I_r=[4096r,\min(4096(r+1),N)),\qquad0\leq r<1038.
\end{equation}
Each ID specifies the full cylinder of all $3^{12}$
possible tails. Different prefixes define disjoint
cylinders and every full digit word has exactly one
prefix, so these cylinders partition all eight trees.
A prefix rejected before depth twelve must still have
an explicit completed record; pruning is not a reason
to omit its ID from the partition.

Separate checks reconstruct every prefix's rank,
digits, and cylinder endpoints and compare them with
both complete decoder output streams. They check
all $N$ IDs and exact endpoints, not a sample or
aggregate length alone. The interval walk checks
that the first endpoint is zero, each start equals
the previous stop, and the last stop is $N$.
Those finite comparisons establish the contents of
the supplied decoder outputs; identifying the outputs
with executions of the intended compiled decoders is
an external provenance assertion, not a Lean theorem.

\section{Sound pruning and exhaustive computation}\label{sec:search}

If $P\subseteq P'$, then $d_P(g)\leq d_{P'}(g)$
for every $g$. Thus a partial set for which some
nonzero difference count exceeds 60 has no
difference-set extension. This is the only
difference-count rejection needed by the search.
At a complete 121-point leaf,
\[
 \sum_{g\ne0}d_P(g)=121\cdot120=242\cdot60.
\]
Consequently all counts being at most 60 is equivalent
to every count being exactly 60. The direct kernel
also checks the exact leaf condition explicitly.

All partial sets are unions of scalar orbits, so
$d_P(61g)=d_P(g)$. Reversing ordered pairs gives
$d_P(-g)=d_P(g)$ as well. There are 36 classes
outside $H$ under $g\mapsto\pm61^ig$, each of
size six, and thirteen nonzero classes inside $H$,
each of size two. Thus there are 49 nonzero
difference classes. The compressed kernel stores
the common count on each class and uses precomputed
option contributions. The direct kernel instead
adds and removes points and maintains the full
elementwise difference vector. Compression is
justified by constancy of contributions on these
classes, rather than by averaging counts.

For every prefix, depth-first search tries all
three digits at each remaining position, unless
the monotone rejection rule applies. An explicit
completion state is required for each prefix;
an absent result, interruption, or partial
continuation is not evidence of exhaustion.
Production searches use the bound~60 and exclude
the relaxed-bound features used for controls.

The two production passes agree on the complete
ordered rows of ID, prefix-validity, node count,
and solution count for all 4,251,528 IDs,
and on node counts in every
one of the 1,038 ranges. With their common node-count
convention, the totals are
\[
\begin{array}{c|r}
\text{depth-twelve prefixes}&4{,}251{,}528\\
\text{search nodes}&108{,}733{,}222{,}341\\
\text{difference-set leaves}&0.
\end{array}
\]
These are traversal counts, not numbers of
inequivalent subsets or an assertion that all
$8\cdot3^{24}$ leaves were visited individually.

A complete additional direct replay reproduced
every range. A separately written third lifting
kernel also exhausted all 1,038 ranges, matching
each range's validity, node, and solution counts
and the digest of its canonical per-prefix row
stream. That kernel has its own recursive
traversal, difference arithmetic, option
construction, ID decoder, and driver. Its
six-cell labeling convention was selected by
testing the two conventions against one range;
this fixes labels only. Permuting the labels
cannot add or remove a subset, and the complete
pass then matched every range.

The compact computational certificate records
each interval and its counts and row-stream
digest. A static check verifies interval geometry
and consistency, while replay recomputes the
rows and compares their counts and digests.
The digest binds a comparison to a particular
row stream; it is not a proof of a universal
nonexistence assertion. The exhaustive executions
and their arithmetic and coverage checks supply
the external computational part of the proof.

\begin{proof}[Proof of Theorem~\ref{thm:main}]
Suppose $D$ were such a difference set.
Lemma~\ref{lem:fixed} gives a fixed translate.
The singleton conditions and complete
classification send its $H$-part legally to
$S_0$. The incidence inversion places its
quotient in the complete 144-array set, and a
legal $S_0$-stabilizing map sends that array
to Table~\ref{tab:cases}. By
Lemma~\ref{lem:lifting}, the resulting
difference set is a leaf of that case. Its
prefix lies in exactly one interval
\eqref{eq:ranges}. Exhaustion rejects every
leaf under that interval by tests that
reject no difference set, a contradiction.
\end{proof}

\section{Completeness checks, controls, and independence}\label{sec:controls}

Agreement of the production kernels alone
does not establish that their shared
decomposition is complete. The proof uses
the full local option-family comparisons,
the separately written exhaustive lifting
pass, the independent finite cover, and
the Lean coverage implication, in addition
to production agreement.

The finite cover checks exact equality of
the singleton sets and legal-map orbit,
the derived quotient set and covered
case orbits, every local option family,
all prefix cylinders, and the range
partition. Three semantic gap controls
exercise these obligations. One deletes
ID~4095 and compensates by duplicating
another ID while preserving the number
of interval records and their summed
length. A second duplicates one legal
option while retaining its cell and
cardinality. A third duplicates a case
of the same orbit size. Each is rejected
with a missing ID, subset, or quotient
image witness, respectively. Their
rejection does not rely on unchanged
file digests or aggregate counts.

Further mutation tests reject altered
scalar partitions, missing maps,
nonimages in the case list, and two
representatives from one quotient
orbit. Replacing a representative
by another member of its own orbit
is accepted, as it should be. An
altered range digest is rejected
by both fresh production kernels.
A changed node count, even with its
aggregate total updated consistently,
is rejected by fresh replay although
it passes the static consistency
check. These tests distinguish
semantic checking from integrity
and formatting checks.

There is a genuine positive discovery
control at parameters $(121,40,13)$
in $C_{121}$: three generic
implementations, without a supplied
witness, recover 88 multiplier-invariant
difference sets at 996,935 nodes.
These are invariant sets without
equivalence reduction. The same
parameters in $C_{11}^2$ give a
negative control; that negative
outcome is not positive discovery
evidence for the present search.

The actual order-243 ternary path
has additional positive controls.
Eight independently constructed
fixed 121-point sets, one in each
case, are supplied through forced
digit paths with relaxed test-only
bound $M=\max_{g\ne0}d_P(g)$.
Both production kernels recover
their exact points at their
specified depth-twelve IDs and
reject bound $M-1$:
\[
\begin{array}{c|r|r}
c&\operatorname{id}&M\\\hline
0&91{,}615&66\\
1&815{,}056&72\\
2&1{,}486{,}729&69\\
3&1{,}750{,}181&72\\
4&2{,}473{,}622&69\\
5&3{,}145{,}295&69\\
6&3{,}214{,}101&69\\
7&3{,}937{,}542&72
\end{array}
\]
For example, the case-zero word is
\[
\texttt{011122200011122200011122}.
\]
It has prefix ID~91,615 and $M=66$.
Each test returns one exact witness
and thirteen tail nodes. These are
bounded-count lifting and witness
controls, not $(243,121,60)$
difference sets: their bounds are
66--72. Neighboring IDs are checked
against their own reconstructed
sets; a neighbor is not assumed
to have no solution at a relaxed
bound. Tests of omitted options,
serialized continuation versus
uninterrupted execution, and
residual symmetries also agree.
Residual symmetries are used for
controls, not for reducing the
production search domain.

A separate 10,000-ID comparison
matches all per-ID rows between
the production kernels and an
additional coordinate-based
implementation. This is a sample
control and does not replace
the complete third-kernel pass
or prove coverage of a range.

The scope of independence is
specific. The compressed and
direct production kernels share
a compilation unit, task and
serialization code, prefix
convention, $S_0$, scalar-partition
and case inputs, cell order,
and option contents. Their
difference representations,
updates, and traversals differ;
they are not independent
derivations of the entire
reduction. The third complete
kernel shares no search or
certificate-driver code with
them, but reads the same
validated scalar-partition
and case files and uses $S_0$.
It reconstructs its own options
and validates orbit membership;
it does not itself discover
the singleton class or quotient
representatives.

Those common mathematical
inputs are instead validated
by separate programs that
regenerate scalar orbits,
enumerate all singleton
candidates, classify quotient
arrays in two different ways,
and check legal full-group
maps and exact case coverage.
Thus the evidence composes
independent input-completeness
checks with a complete separate
lifting exhaustion. A complete
search fed directly by a newly
self-generated chain of all
front-end data is not claimed.
The common mathematical
implications are justified by
the arguments and formalization,
not inferred from matching
search totals.

\section{The exact formalization boundary}\label{sec:formal}

The Lean definition of the
group is the product of the
three residue groups of
moduli $3,9,9$. Its
\texttt{IsDiffSet} predicate
is exactly cardinality~121
and ordered-difference
count~60 at every nonzero
element. The formalization
proves the first-multiplier
specialization by the
classical group-ring argument,
as well as centering, the
index-three moments, the
singleton congruences,
legal normalizations,
quotient inversion, lifting,
pruning, and ID partition.

In particular,
\texttt{coverage} proves the
following implication, with
no hypotheses about outputs
of a finite enumerator or
of a search: if $D$ is a
difference set, then there
exist $u\in G$, a word $w$
of legal affine maps, an
integer $r<1038$, an
$\operatorname{id}\in I_r$,
and a twelve-digit tail
$\boldsymbol t$ such that
\begin{equation}\label{eq:coverage}
 F_w(D-u)=\mathcal L\left(
 \left\lfloor\frac{\operatorname{id}}{531441}\right\rfloor,
 \operatorname{digits}(\operatorname{id},\boldsymbol t)
 \right).
\end{equation}
Here the first twelve digits
are the leading-zero-preserving
base-three expansion of
$\operatorname{id}\bmod531441$
and the last twelve are
$\boldsymbol t$. The map
$F_w$ is an additive
automorphism followed by a
translation in $H$.
This is an existence
statement; the normalizing
map or normalized
representative for the
original $D$ need not be
unique. The interval of
each ID is unique.

The formal singleton proof
uses its own finite
classification: fourteen
layer-zero representatives,
64 layer-one shear classes,
52,220 patterns, and twelve
survivors with legal words
to $S_0$. These numbers
describe that classification,
not the 5,335,092-mask
enumeration in
\eqref{eq:maskcount}.
The quotient proof checks
all $2^{13}$ assignments
and the legal witnesses
to the eight actual
case rows. Generated
witness tables have their
required properties and
completeness checked in
the kernel; they are not
imported as axioms.

For an interval $I_r$,
define the semantic
assertion
\begin{equation}\label{eq:clear}
\begin{split}
 \operatorname{RangeClear}(r):\quad
 &\text{for every }\operatorname{id}\in I_r
 \text{ and every twelve-digit tail }\boldsymbol t,\\
 &\mathcal L\left(
 \left\lfloor\frac{\operatorname{id}}{531441}\right\rfloor,
 \operatorname{digits}(\operatorname{id},\boldsymbol t)
 \right)\text{ is not a difference set.}
\end{split}
\end{equation}
The final Lean theorem
\texttt{no\_difference\_set}
is conditional on the family
\[
 \forall r<1038,\qquad
 \operatorname{RangeClear}(r).
\]
These are 1,038 semantic
assertions. A receipt
formulation packages each
assertion as the
\texttt{clear} field of
$\operatorname{Receipt}(r,\operatorname{digest}_r)$;
the digest is only a
label, not additional
mathematical evidence.
Lean does not compute
SHA-256 here, interpret
C++ output, or derive
the assertions from
digest strings. The external
exhaustive computations
in Section~\ref{sec:search}
supply that part of the
computer-assisted proof.

Equality of the Lean
leaf table with the
implemented search
options is also an
external comparison,
covering all 576 options
and thirteen singleton
points, strengthened by
the complete compiled
option dumps and separate
definition-based checks.
It is neither an omitted
hypothesis of
\texttt{coverage} nor a
formal proof of the C++
implementation, compiler,
or execution environment.
The finite cover's
decoder comparisons
likewise retain the
execution-provenance
boundary stated in
Section~\ref{sec:lifting}.

The recorded kernel
checks of the reductions
and conditional final
theorem use only
\texttt{propext},
\texttt{Classical.choice},
and \texttt{Quot.sound}.
Planted formal errors
are rejected: dropping
the last range, pruning
when a count reaches
$\lambda$ rather than
exceeds it, weakening
the exact leaf criterion
to $\lambda+1$, and
omitting an admissible
singleton survivor.
The first three false
statements also have
formal counterexamples.
These checks establish
the stated reduction
scope; they do not
turn the external
search into a
whole-theorem Lean
proof from the
definitions.

\begin{thebibliography}{99}

\bibitem{ArasuMa2001}
K. T. Arasu and S. L. Ma,
\emph{A nonexistence result on difference sets, partial difference sets
and divisible difference sets},
Journal of Statistical Planning and Inference \textbf{95} (2001), 67--73.
\href{https://doi.org/10.1016/S0378-3758(00)00278-0}{doi:10.1016/S0378-3758(00)00278-0}.

\bibitem{ChenXiangSehgal1994}
Y. Q. Chen, Q. Xiang, and S. K. Sehgal,
\emph{An exponent bound on skew Hadamard abelian difference sets},
Designs, Codes and Cryptography \textbf{4} (1994), 313--317.
\href{https://doi.org/10.1007/BF01388647}{doi:10.1007/BF01388647}.

\bibitem{Gordon2019}
Daniel M. Gordon,
\emph{The La Jolla Difference Set Repository},
presentation, August~3, 2019, slides~11 and~38
(slide~38 is physical PDF page~49).
\url{https://cargo.wlu.ca/ArasuFest/ArasuFest_talks/Dan_Gordon.pdf}.

\bibitem{GordonData}
Daniel M. Gordon,
\emph{La Jolla Difference Sets Repository},
archived dataset, version~1.3, April~24, 2026.
\href{https://doi.org/10.5281/zenodo.19740740}{doi:10.5281/zenodo.19740740}.
Repository and definition of entries:
\url{https://dmgordon.org/difference-sets/}.

\bibitem{Hall1956}
Marshall Hall Jr.,
\emph{A Survey of Difference Sets},
Proceedings of the American Mathematical Society
\textbf{7} (1956), 975--986.
\href{https://doi.org/10.1090/S0002-9939-1956-0082502-7}{doi:10.1090/S0002-9939-1956-0082502-7}.

\bibitem{HeWu2026}
Zhiwen He and Jiayan Wu,
\emph{Nonexistence results and constructions for signed difference sets},
preprint, arXiv:2608.22936v2, September~20, 2026
(first version August~24, 2026).
\url{https://arxiv.org/abs/2608.22936v2}.

\bibitem{Lander1983}
E. S. Lander,
\emph{Symmetric Designs: An Algebraic Approach},
London Mathematical Society Lecture Note Series, vol.~74,
Cambridge University Press, Cambridge, 1983.
\href{https://doi.org/10.1017/CBO9780511662164}{doi:10.1017/CBO9780511662164}.

\bibitem{Lehmer1953}
Emma Lehmer,
\emph{On residue difference sets},
Canadian Journal of Mathematics \textbf{5} (1953), 425--432.
\href{https://doi.org/10.4153/CJM-1953-047-3}{doi:10.4153/CJM-1953-047-3}.

\bibitem{LeungMaSchmidt2004}
K. H. Leung, S. L. Ma, and B. Schmidt,
\emph{Nonexistence of abelian difference sets: Lander's conjecture
for prime power orders},
Transactions of the American Mathematical Society
\textbf{356} (2004), 4343--4358.
\href{https://doi.org/10.1090/S0002-9947-03-03365-8}{doi:10.1090/S0002-9947-03-03365-8}.

\bibitem{LeungSchmidt2005}
K. H. Leung and B. Schmidt,
\emph{The Field Descent Method},
Designs, Codes and Cryptography \textbf{36} (2005), 171--188.
\href{https://doi.org/10.1007/s10623-004-1703-7}{doi:10.1007/s10623-004-1703-7}.

\bibitem{Miller2025}
Niklas Miller,
\emph{Forbidden multipliers in abelian difference sets},
preprint, arXiv:2511.10231v1, November~13, 2025.
\url{https://arxiv.org/abs/2511.10231v1}.

\bibitem{Paley1933}
R. E. A. C. Paley,
\emph{On Orthogonal Matrices},
Journal of Mathematics and Physics \textbf{12} (1933), 311--320.
\href{https://doi.org/10.1002/sapm1933121311}{doi:10.1002/sapm1933121311}.

\bibitem{Turyn1965}
R. J. Turyn,
\emph{Character sums and difference sets},
Pacific Journal of Mathematics \textbf{15} (1965), 319--346.
\href{https://doi.org/10.2140/pjm.1965.15.319}{doi:10.2140/pjm.1965.15.319}.

\bibitem{VeraGarcia1997}
A. Vera L\'opez and M. A. Garc\'ia S\'anchez,
\emph{On the existence of abelian difference sets with $100<k\leq150$},
Journal of Combinatorial Mathematics and Combinatorial Computing
\textbf{23} (1997), 97--112; Table~1, p.~105.
\href{https://combinatorialpress.com/jcmcc-articles/volume-023/on-the-existence-of-abelian-difference-sets-with-100-k-leq-150/}{Publisher's article record}.

\end{thebibliography}

\end{document}
